% FailWhere: Can Vision-Language Models Self-Localize Their Own Failures?
% DMLR submission version
%
% To compile: requires dmlr2e.sty (in this directory, from
% https://github.com/JmlrOrg/dmlr-style-file).
% Then: pdflatex paper_dmlr && bibtex paper_dmlr && pdflatex paper_dmlr && pdflatex paper_dmlr
%
% DMLR is single-blind: author info, real repo URL, and affiliations are shown.

\documentclass[twoside,11pt]{article}

\usepackage[preprint]{dmlr2e}
\usepackage{booktabs}
\usepackage{multirow}
\usepackage{amsmath}
\usepackage{url}

% Short headings: running head and authors last names
\ShortHeadings{FailWhere: VLM Failure Self-Localization}{Wandile}
\firstpageno{1}
% OpenReview link shown on the first page; update with the forum URL after submission.
\def\openreview{\url{https://openreview.net/group?id=DMLR}}

\begin{document}

\title{FailWhere: Can Vision-Language Models Self-Localize Their Own Failures?}

\author{\name Subodh Wandile \email subodh.anant.wandile@toasttab.com \\
       \addr Toast, Inc.\\
       Boston, MA, USA}

\maketitle
\begin{abstract}
When a Vision-Language Model (VLM) answers a visual question incorrectly, can it
accurately identify which spatial region of the image caused that error?
We introduce \textsc{FailWhere}, a benchmark that tests this capacity directly:
for each answer (correct or wrong), we ask the model to name the 3$\times$3 grid
region most responsible using a structured \textsc{region:} tag.
On controlled synthetic images (500 samples, Claude Haiku), structured
self-localization accuracy is only 9\% against an answer-cell ground truth,
despite 75\% keyword-level mention — a verbal-commitment gap showing the model
hedges across regions in text but cannot commit.
On real COCO photographs (VSR dataset, 142 failures), four converging lines of
evidence show the \textsc{region:} tag is generic image re-reading rather than
failure self-knowledge:
\textbf{(1)}~hit rates on failures match correct answers within 1.1\,pp
(\textsc{tost} equivalence confirmed) — correctness-invariant spatial awareness;
\textbf{(2)}~a different model (Sonnet-5-5) that never produced the wrong answers
attributes Haiku's failures equally well ($+$25.4\,pp above null, $p{=}0.002$) —
no privileged access to the answering model's state;
\textbf{(3)}~center tags carry zero causal signal; and
\textbf{(4)}~non-center tags land $+$23--25\,pp above causal ground truth across
models.
Verbal self-attribution is prior-dominated in the modal case (${\sim}54\%$ center)
but meaningfully informative when the model commits to a specific non-center region.
We release the benchmark and all code.
\end{abstract}

\begin{keywords}
  vision-language models, interpretability, failure attribution,
  benchmarks, spatial reasoning, faithfulness
\end{keywords}



%% ===================================================================
\section{Introduction}
%% ===================================================================

Modern VLMs achieve strong performance on visual question answering benchmarks,
yet failures remain common on spatial reasoning tasks.
Fixing these failures requires understanding \emph{why} a model failed —
which region of the image triggered the wrong answer.
This is an open problem: prior attribution work either uses model internals
(attention maps, gradients) or focuses on cases where the model was correct.

\paragraph{The gap.}
We map prior work onto a 2$\times$2 table
(Table~\ref{tab:2x2}) across two dimensions:
(1)~correct or wrong answer, and
(2)~verbal (black-box) or internal (white-box) attribution.
\textsc{FailWhere} occupies the empty cell: verbal, black-box, failure-case localization.

\paragraph{Contributions.}
\textbf{(1)}~We formalize \emph{VLM failure self-localization}: given a failure,
can the model name the responsible image region using only its verbal output?
\textbf{(2)}~We introduce \textsc{FailWhere}: a synthetic benchmark with controlled
3$\times$3 grid images and counterfactual-masking ground truth (CF-GT).
\textbf{(3)}~We provide the first empirical results:
9\% self-localization on synthetic images;
a two-regime structure on real COCO photos;
and a Sonnet-vs-Haiku control that resolves the ``self'' question cleanly.
\textbf{(4)}~We release benchmark code and all result files.%
\footnote{Code, data, and result JSONs:
\url{https://github.com/subodhwandile/vlm-failure-bench}
(MIT license; maintained via the repository issue tracker).}


%% ===================================================================
\section{Related Work}
%% ===================================================================

\begin{table}[t]
\centering
\small
\begin{tabular}{lll}
\toprule
 & Correct answers & \textbf{Wrong answers} \\
\midrule
\textbf{Verbal} & CRAFT \citep{parcalabescu-frank-2023-craft} & \textsc{FailWhere} (ours) \\
\textbf{Internal} & EDCT-Bench \citep{zang-2026-edct}; GLIMPSE & MLLMs-Know \citep{zhang-2025-mllms} \\
\bottomrule
\end{tabular}
\caption{2$\times$2 positioning of FailWhere among prior attribution work.
The verbal$\times$failure cell was previously empty.}
\label{tab:2x2}
\end{table}

\paragraph{White-box failure attribution.}
\citet{zhang-2025-mllms} use attention maps and gradient overlays to identify
which image regions VLMs attend to when making errors.
\textsc{FailWhere} asks whether verbal self-reports approximate this signal without
model internals.

\paragraph{Verbal correct attribution.}
CRAFT \citep{parcalabescu-frank-2023-craft} tests whether VLMs explain their
\emph{correct} answers faithfully.
\textsc{FailWhere} asks the same question on failures, where no correct explanation
is available to copy from.

\paragraph{VLM self-knowledge and CoT faithfulness.}
\citet{madsen-etal-2024-faithfulness} show that chain-of-thought explanations are
often unfaithful post-hoc rationalizations.
Our verbal-commitment gap (75\% keyword vs.\ 9\% structured accuracy) is a
spatial-domain instance of this finding.
\citet{kadavath-2022-knowwhat} show LLMs can self-calibrate factual confidence but
struggle with spatial confidence.
\citet{shen-2023-visual-cot} find VLM spatial descriptions less faithful than
factual ones.
Seeing-but-Not-Believing \citep{seeing-2025} shows VLMs can encode a scene
correctly yet fail to believe their own percept — directly motivating our question.
CLIP Cross-Eyed \citep{clipcross-2026} finds that verbalized spatial descriptions
misalign with internal spatial representations, consistent with our center-prior
finding: the verbal channel reports a distributional prior, not an introspective readout.


%% ===================================================================
\section{The \textsc{FailWhere} Benchmark}
%% ===================================================================

\subsection{Synthetic Benchmark}

Images are 300$\times$300\,px, divided into a uniform 3$\times$3 grid.
Each cell contains one object (shape $\times$ color).
Two question types:

\begin{itemize}
  \item \textbf{Adjacent-distractor}: same shape in target and adjacent cell,
    different colors; model answers with the distractor color.
  \item \textbf{Relative-position}: ``What is [dir] of the [ref-object]?'' —
    requires two-step spatial reasoning; model misidentifies reference or shifts
    one cell.
\end{itemize}

\noindent Ground truth: \emph{answer-cell GT} (cell containing the correct answer).

\subsection{Real-Image Extension: VSR}

We use the Visual Spatial Reasoning dataset \citep{liu-2023-vsr}
(VSR, \texttt{cambridgeltl/vsr\_random}, split \texttt{test}),
which pairs COCO train2017 photographs with 66 spatial-relation caption types.
Task: classify each caption as True or False.
We use the first 500 items, at temperature~$= 0.0$.

\subsection{Counterfactual Masking Ground Truth (CF-GT)}

For each failure, we mask each of 9 cells with solid gray (128,128,128)
and re-run the VLM.
The cell whose masking corrects the wrong answer is CF-GT.
If no cell fixes it: \texttt{none\_found}.
If several cells each fix it: \texttt{multiple}.
Only single-cell \emph{clean} outcomes ($n{=}22/142$, 15.5\%) are used to score
tag accuracy; this 15.5\% rate is the hard attribution ceiling for this dataset.

\subsection{Self-Localization Protocol}

Three prompt conditions (full templates in Appendix~\ref{app:prompts}):

\begin{itemize}
  \item \textbf{Condition A (Neutral)}: ``Which region most influenced your answer?''
    — all items.
  \item \textbf{Condition B (Stateless)}: ``Your answer was wrong. Which region
    caused it?'' — failures only.
  \item \textbf{Condition C (Multi-turn)}: original T/F exchange prepended as
    context before attribution prompt — failures only.
\end{itemize}

\noindent The \textsc{region:} tag uses exactly one of nine labels
(top-left, top-center, \ldots, bottom-right).
Parse rate: 100\% on the VSR run.


%% ===================================================================
\section{Results}
%% ===================================================================

\subsection{Synthetic Benchmark}
\label{sec:synthetic}

\begin{table}[t]
\centering
\small
\begin{tabular}{lr}
\toprule
Metric & Value \\
\midrule
Samples & 500 \\
Error rate & 8.8\% (44/500) \\
Structured tag accuracy (answer-cell GT) & \textbf{9\%} (4/44) \\
Keyword match (region mentioned in text) & 75\% (33/44) \\
Mean IoU (structured tag) & 0.091 \\
Failures by type (relative-pos / adj-distract) & 40 / 4 \\
\bottomrule
\end{tabular}
\caption{Synthetic benchmark results (Claude Haiku, $n{=}500$, temp$=0$).}
\label{tab:synthetic}
\end{table}

Table~\ref{tab:synthetic} shows a \emph{verbal-commitment gap}: 75\% keyword
mention vs.\ 9\% structured accuracy.
The model mentions the correct region in its explanation but commits to a different
one in the \textsc{region:} tag — hedging across many regions in free text while
defaulting to a wrong heuristic when forced to name one.
Both the 200-sample pilot (0\% / 17 failures) and this 500-sample run confirm
near-zero self-localization under answer-cell GT.

\subsection{Real Images: Center Bias}
\label{sec:center}

On 142 failures (VSR, temp$=0$), the model tags ``center'' in 45.8\% of
Condition~B responses — 4.1$\times$ the uniform prior (11.1\%).
The per-instance COCO bbox analysis (Section~\ref{sec:coco}) shows this reflects
veridical image statistics: COCO subjects are in the center cell 49.3\% of the time.
The center-bias is not an error artifact — it tracks object distributions.

\subsection{Correctness-Invariant Spatial Awareness}
\label{sec:coco}

\begin{table}[t]
\centering
\small
\begin{tabular}{lrrr}
\toprule
 & Correct ($n{=}336$) & Failed ($n{=}131$) & $\Delta$ \\
\midrule
Subject cell hit & 46.1\% & 45.0\% & $-1.1$\,pp \\
Object cell hit & 26.2\% & 28.2\% & $+2.1$\,pp \\
Either cell hit & 59.2\% & 61.8\% & $+2.6$\,pp \\
\midrule
\emph{Baseline: uniform} & 11.1\% & 11.1\% & — \\
\emph{Baseline: perm-null} & 23.2\% & 22.9\% & — \\
\emph{Baseline: always-center} & 49.1\% & 49.6\% & — \\
\bottomrule
\end{tabular}
\caption{Per-instance COCO bbox hit rates for neutral tag (Cond.\ A),
comparing correct vs.\ failed answers ($n{=}467$ scored items, 93.4\% join rate).
\textsc{tost} equivalence test ($\pm 8$\,pp bounds): $p{=}0.86$ (subject),
$p{=}0.62$ (either) — cannot reject equivalence.}
\label{tab:coco}
\end{table}

Table~\ref{tab:coco} shows three decisive findings:

\textbf{(1)}~Hit rates on failed answers match correct answers within 2.6\,pp
(all within 95\% Wilson CIs).
TOST equivalence ($\pm 8$\,pp) is confirmed ($p{\geq}0.62$):
the \textsc{region:} tag does not distinguish whether the answer was correct.

\textbf{(2)}~The verbal report is prior-dominated:
``always predict center'' achieves 49.1--49.6\% subject-hit, \emph{beating} the
model's 46.1--45.0\%.
The \textsc{region:} tag captures image-object priors but adds no discriminative
information beyond them.

\textbf{(3)}~Center sensitivity analysis:
when the subject IS in center, the model tags center 57.6\% (correct) /
53.8\% (failed); when it is not, the model still tags center 25.1\% / 31.8\%.
As a binary center-detector: sensitivity ${\approx}56\%$, FPR ${\approx}28\%$,
AUC ${\approx}0.64$ — a weak detector.

\subsection{Counterfactual Attribution: Two-Regime Structure}
\label{sec:cf}

\begin{table}[t]
\centering
\small
\begin{tabular}{lrrr}
\toprule
 & Neutral (A) & Stateless (B) & Multi-turn (C) \\
\midrule
Clean hit rate ($n{=}22$) & 27.3\% & 31.8\% & 36.4\% \\
Perm-null & 13.5\% & 14.0\% & 13.7\% \\
Excess over null & $+$13.8\,pp & $+$17.8\,pp & $+$22.7\,pp \\
\midrule
\emph{Tag = center} ($n{=}3$) & 33.3\% & 66.7\% & 33.3\% \\
excess & $+$0.0\,pp & $+$0.0\,pp & $+$0.0\,pp \\
\emph{Tag $\neq$ center} ($n{=}19$) & 26.3\% & 26.3\% & 26.3\% \\
excess & $+$14.9\,pp & $+$16.4\,pp & $+$16.4\,pp \\
\bottomrule
\end{tabular}
\caption{Tag vs.\ CF-GT accuracy on clean single-cell subset ($n{=}22$).
Center excess is exactly null across all conditions; non-center excess is
stable at $+$15--16\,pp — identical across conditions.}
\label{tab:cf}
\end{table}

\paragraph{Two-regime model.}
Table~\ref{tab:cf} reveals a clean bifurcation:

\textbf{Center tags carry zero causal signal.}
The center hit rate equals the permuted-null and the always-center baseline
(both ${\approx}14\%$); excess is 0\,pp across all three conditions.
Center self-attribution provides no information about what actually caused the error.

\textbf{Non-center tags carry real signal.}
When the model escapes the center prior, it hits CF-GT at 26.3\% vs.\ a
17.9\% null — a consistent $+$16\,pp excess across all three conditions,
approximately $2.3\times$ above chance.

\textbf{Context (Condition~C) does not improve attribution.}
The non-center hit rate is identical across all three conditions (26.3\%),
confirming the two-regime structure is prompt-invariant.
The multi-turn gain on the overall number ($+$22.7\,pp) is a center-stratum
artifact ($n{=}3$, noise).

\paragraph{Failure typology.}
Of 142 CF runs: \texttt{none\_found} 51.4\%, \texttt{multiple} 33.1\%,
clean 15.5\%.
Most VSR failures are causally diffuse at 3$\times$3 granularity — no single
cell explains the error.
The clean ceiling is the paper's hardest constraint: any tag-accuracy claim on
this dataset is bounded by 15.5\%, not 100\%.

\paragraph{Label asymmetry.}
True-caption failures (model said False, $n{=}104$): 60.6\% none\_found,
11.5\% clean — holistic spatial-comprehension failures.
False-caption failures (model said True, $n{=}38$): 26.3\% none\_found,
26.3\% clean — more localized, region-specific evidence.

\subsection{Other-Model Baseline: Sonnet Reads Haiku's Failures}
\label{sec:sonnet}

\begin{table}[t]
\centering
\small
\begin{tabular}{lrrrr}
\toprule
 & \multicolumn{2}{c}{Sonnet-5-5} & \multicolumn{2}{c}{Haiku-4-5} \\
 & Neutral & Stateless & Neutral & Stateless \\
\midrule
Clean hit ($n{=}22$) & 31.8\% & \textbf{36.4\%} & 27.3\% & 31.8\% \\
Excess over null & $+$20.4\,pp & $+$25.4\,pp & — & — \\
$p$-value (one-sided) & 0.007 & \textbf{0.002} & — & — \\
Center bias & 39.7\% & 41.4\% & 41.8\% & 45.4\% \\
\bottomrule
\end{tabular}
\caption{Sonnet-5-5 attributing Haiku-4-5's 142 VSR failures vs.\ Haiku's
own attributions, scored against Haiku's CF-GT.
Sonnet matches or slightly exceeds Haiku despite never having produced the
wrong answers. Sonnet stateless result is the paper's strongest significant
number: $+$25.4\,pp, $p{=}0.002$.}
\label{tab:sonnet}
\end{table}

\paragraph{Design.}
Self-knowledge control: Sonnet-5-5 is given the image, caption, and
Haiku's wrong answer, then produces a \textsc{region:} tag.
If Sonnet matches Haiku's attribution accuracy, the tag reflects generic image
re-reading; if it is worse, self-knowledge is required.

\paragraph{Result.}
Table~\ref{tab:sonnet} shows Sonnet matches or slightly exceeds Haiku:
31.8--36.4\% vs.\ 27.3--31.8\%.
The Sonnet stateless result is the paper's strongest significant number:
$+$25.4\,pp above null, $p{=}0.002$.
A model that never produced the wrong answer attributes the failure as well as
the model that did.

\paragraph{Implication.}
\textbf{The \textsc{region:} tag is a generic re-reading of the image, not
privileged access to the answering model's internal state.}
Combined with Section~\ref{sec:coco}'s correctness-invariance, Sections
\ref{sec:cf}--\ref{sec:sonnet} jointly establish the four-pillar story of
Section~\ref{sec:discussion}.


%% ===================================================================
\section{Discussion}
\label{sec:discussion}
%% ===================================================================

\paragraph{Four pillars.}
The paper's results converge on a complete structured negative result:

\begin{enumerate}
  \item \textbf{Reads the image, not the mind} — Sonnet attributes Haiku's failures
    equally well ($+$25.4\,pp, $p{=}0.002$).
    No privileged model-state access is needed or demonstrated.
  \item \textbf{Doesn't know if the answer was wrong} — identical hit rates on
    correct and failed answers ($\Delta = -1.1$\,pp, TOST $\pm 8$\,pp confirmed).
  \item \textbf{Commits reasonably when it escapes center} — non-center tags land
    $+$15--25\,pp above null across conditions and models.
  \item \textbf{Most reports never escape} — ${\sim}54$\% are center-tagged;
    center tags carry zero causal signal.
\end{enumerate}

\paragraph{The center bias is not spurious.}
COCO subjects sit in the center cell 49.3\% of the time; the model's 42\%
neutral-center rate reflects real image statistics, not a failure artifact.
However, the center prior is counterproductive under failure conditions:
causal cells are 86\% non-central (19/22 clean), so always-center achieves
only 13.6\% CF-GT accuracy vs.\ the model's 27--32\%.
The center prior is informative for object localization and harmful for
failure attribution — the two tasks are confounded only at the level of the
verbal tag.

\paragraph{Two GT definitions and their ceiling.}
\textit{Answer-cell GT} (cell containing the correct answer) targets capability;
\textit{CF-GT} (cell whose masking fixes the error) targets causality.
On the VSR dataset, 84.5\% of failures have no clean CF-GT ($\geq 2$ cells or
none required); the effective attribution ceiling is 15.5\%.
Any result reporting tag-vs-CF-GT accuracy must be interpreted relative to this
ceiling, not 100\%.

\paragraph{Limitations.}
(1)~Caption-to-COCO bbox matching uses heuristic noun extraction (93.4\% join rate).
(2)~All experiments use Claude Haiku and Sonnet; cross-family replication
(GPT-4V, Gemini) remains future work.
(3)~Task scope is binary T/F (VSR); generalization to open-ended VQA is untested.
(4)~The stateless prompt asks for post-hoc attribution; the model has no memory
of its visual processing chain.
(5)~CF masking at 3$\times$3 has a resolution limit: synthetic adjacent-distractor
cases return \texttt{multiple} (4/4) even when the causal cell is known by
construction.
(6)~Fix-rate secondary analysis ($n{=}64$, $+$4.9\,pp, $p{=}0.20$) is
directionally consistent with the binary result but underpowered.


%% ===================================================================
\section{Conclusion}
%% ===================================================================

We introduce \textsc{FailWhere} and show that VLM verbal self-localization is
generic image re-reading with a center prior — not failure self-knowledge.
The result is a complete structured negative: the \textsc{region:} tag
(1)~requires no model-state access (Sonnet control), (2)~does not detect
failure (correctness-invariance), (3)~is informative when non-center
($+$15--25\,pp), and (4)~is mostly captured by the center prior (${\sim}54\%$
of reports).
These findings motivate both better introspection training objectives and
more precise evaluation frameworks that distinguish self-report from genuine
failure attribution.


%% ===================================================================
%% ===================================================================
\impact{This work studies whether vision-language models can accurately
identify the image region responsible for their own errors.
Positively, a reliable self-localization capacity would support error
diagnosis, targeted model debugging, and more trustworthy deployment of
VLMs in spatial-reasoning settings (e.g., robotics, document
understanding, assistive tools). Our findings are cautionary in a way
that also benefits practitioners: a fluent, confident \textsc{region:}
tag is largely generic image re-reading driven by a center prior, so
downstream systems should not consume verbal failure attributions as
ground-truth causal evidence. Negatively, the same finding could
encourage over-trust in models that \emph{appear} self-aware;
misattributed explanations could misdirect debugging effort or mask
systematic failure modes. The benchmark uses public datasets (COCO
images via the VSR dataset) and synthetic images; it contains no
personal data and poses no privacy concerns.}

\acks{The author thanks the creators of the Visual Spatial Reasoning
dataset and the COCO dataset for public data access, and Anthropic for
API access to the evaluated models.
This research received no external funding and was conducted in a
personal capacity; the author declares no competing interests beyond
their employment at Toast, Inc., which had no role in this work.}

\vskip 0.2in
\bibliography{refs}
% dmlr2e sets the natbib bibliography style (plainnat) automatically.

\appendix

\section{Prompt Templates}
\label{app:prompts}

All prompts use temperature $= 0.0$ and model
\texttt{claude-haiku-4-5-20251022} (public Anthropic API) /
\texttt{global.anthropic.claude-haiku-4-5-20251001-v1:0} (Bedrock).
Sonnet baseline (\S\ref{sec:sonnet}): \texttt{claude-sonnet-5-5} /
\texttt{global.anthropic.claude-sonnet-5-5};
\texttt{temperature} omitted on Bedrock (deprecated parameter for this model).

\subsection{True/False Question (all conditions)}
\begin{verbatim}
Is the following statement true or false?

Statement: "{caption}"

Answer with only "True" or "False".
\end{verbatim}

\subsection{Condition A — Neutral Self-Localization}
\begin{verbatim}
You just answered the following question about this image:
"{caption}"

Your answer was: "{model_answer}"

Which region of the image most influenced your answer?

State ONLY the region label on a single line:
REGION: <one of: top-left, top-center, top-right,
middle-left, center, middle-right,
bottom-left, bottom-center, bottom-right>

Then briefly explain why that region influenced your answer.
\end{verbatim}

\subsection{Condition B — Stateless Failure-Specific}
\begin{verbatim}
You answered the following question INCORRECTLY.

Question: Is the following statement true or false?
Statement: "{caption}"

Your answer was: "{model_answer}" (this was wrong).

Which region of the image most influenced your incorrect answer?

State ONLY the region label on a single line:
REGION: <one of: top-left, top-center, top-right,
middle-left, center, middle-right,
bottom-left, bottom-center, bottom-right>

Then explain:
1. What visual content in that region led to the wrong answer
2. Why that region caused your error
\end{verbatim}

\subsection{Condition C — Multi-Turn}
\begin{verbatim}
[Turn 1 — user]
Is the following statement true or false?
Statement: "{caption}"
Answer with only "True" or "False".

[Turn 1 — assistant]
{model_answer}

[Turn 2 — user]
That answer was incorrect. Looking back at the image,
which region most influenced your incorrect answer?

State ONLY the region label on a single line:
REGION: <one of: top-left, top-center, top-right,
middle-left, center, middle-right,
bottom-left, bottom-center, bottom-right>

Then explain:
1. What visual content in that region led to the wrong answer
2. Why that region caused your error
\end{verbatim}

\subsection{Grid Cell Definitions}

The image is divided into a uniform 3$\times$3 grid:

\begin{center}
\small
\begin{tabular}{lrrrr}
\toprule
Cell & $x_0$ & $y_0$ & $x_1$ & $y_1$ \\
\midrule
top-left    & 0   & 0   & 1/3 & 1/3 \\
top-center  & 1/3 & 0   & 2/3 & 1/3 \\
top-right   & 2/3 & 0   & 1   & 1/3 \\
middle-left & 0   & 1/3 & 1/3 & 2/3 \\
center      & 1/3 & 1/3 & 2/3 & 2/3 \\
middle-right& 2/3 & 1/3 & 1   & 2/3 \\
bottom-left & 0   & 2/3 & 1/3 & 1   \\
bottom-center&1/3 & 2/3 & 2/3 & 1   \\
bottom-right& 2/3 & 2/3 & 1   & 1   \\
\bottomrule
\end{tabular}
\end{center}

\noindent CF masking: cells are filled with solid gray (128, 128, 128).
Images encoded as PNG base64 in API requests to avoid JPEG re-compression
artifacts in masked regions.

\end{document}
